\section{Evaluated full-matrix traces in the finite-boundary case}
\label{sec:finite-boundary-trace-corollary}

This corollary combines two separately checked operator
constructions. MA-MO-027 extends the original full matrix
under its finite-boundary and uniform-interior-gap hypotheses.
MA-MO-028 evaluates the actual ideal or positive-source
endomorphism at every width. Their combination removes
the reserved ideal-trace evaluation in (MA-MO.704).
It does not weaken the full-source hypotheses of MA-MO-027.

Retain precisely (MA-MO.700)--(MA-MO.702), including
the strict inequality \(s>2m-1\) at every boundary
zero of multiplicity \(m\). Use the structural nontrivial
character setting of MA-MO-022, or the rational principal
setting \(S=\{\infty\}\) with the actual polar completion.
The cutoff is the same arbitrary real smooth step.
For a compact smooth test put
\[
 \begin{aligned}
 G(z)&=\mathcal Mg(z),&
 \psi&=1-\phi-\phi^\vee,\\
 Z_\omega(g)&=\sum_\rho m_\omega(\rho)G(\rho-1/2),&
 p(g)&=G(-1/2)+G(1/2),\\
 \tau_{\omega,S}(g,\phi)
   &=\frac1{2\pi}\int_{\mathbb R}
             G(it)\frac{m_\omega'(it)}{m_\omega(it)}\,dt
                         +g(0)\int_{\mathbb R}\psi(x)\,dx .
 \end{aligned}
 \tag{MA-MO.850}
\]
The local multiplier is the actual one, including all
finite Fourier phases and conductor/different monomials.
For the rational case write \(Z_0,\tau_0\).
The completed-zero sum is absolutely convergent.

\paragraph{MA-MO-029 (finite-boundary full traces and genuine determinants).}
In the nontrivial character setting, the original full
extension from MA-MO-027 satisfies
\[
 \operatorname{Tr}_H\mathscr D_{\eta,g}
                 =\tau_{\eta,S}(g,\phi)-Z_\eta(g).
 \tag{MA-MO.851}
\]
Its actual ideal restriction is the all-width operator
\(\mathcal B_{\eta,g}\) of MA-MO-028, and its determinant is
\[
 \begin{aligned}
 \det_H(1-z\mathscr D_{\eta,g})
   &=\det_{J_\eta}(1-z\mathcal B_{\eta,g})
                                \det_{Q_\eta}(1+zC_g)\\
   &=\det_{J_\eta}(1-z\mathcal B_{\eta,g})
       \prod_\rho\bigl(1+zG(\rho-1/2)\bigr)^{m_\eta(\rho)} .
 \end{aligned}
 \tag{MA-MO.852}
\]
Both equalities hold for every complex \(z\).
For \(g(0)=0\) the trace has the exact finite-prime value
\[
 \begin{aligned}
 \operatorname{Tr}_H\mathscr D_{\eta,g}
  =\sum_{\substack{\mathfrak p\notin S\\k\ge1}}
    &\frac{\log N\mathfrak p}{(N\mathfrak p)^{k/2}}
       \bigl[
        c_{\mathfrak p}^{\,k}g(k\log N\mathfrak p)\\
    &\hspace{9em}
       +\overline{c_{\mathfrak p}}^{\,k}g(-k\log N\mathfrak p)
       \bigr].
 \end{aligned}
 \tag{MA-MO.853}
\]
Here \(c_{\mathfrak p}\) is the primitive unramified
character value. No contributing prime is implicitly
added to \(S\).

In the rational setting put
\[
 \begin{gathered}
 \mathcal W=H\oplus\mathbb C^2,\qquad
 \mathcal W_+=J\oplus\mathbb C^2,\qquad Q=H/J,\\
 L_g=\jmath_+\mathcal A_{+,g}\in\mathcal S_1(\mathcal W_+),\\
 F_g=C_g(1-P_-),\qquad
 \mathscr E_g=F_g-\mathscr D_g .
 \end{gathered}
 \tag{MA-MO.854}
\]
The return \(\mathcal A_{+,g}\) is now the all-width
one in (MA-MO.829). The operator \(\mathscr E_g\)
extends the original \(C_g(1-P_+)\), preserves
\(\mathcal W_+\), restricts there to \(-L_g\),
and induces \(C_g\) on \(Q\). In particular
\[
 \begin{aligned}
 \operatorname{Tr}_{\mathcal W}\mathscr E_g
                       &=-\tau_0(g,\phi)+Z_0(g),\\
 \operatorname{Tr}_{\mathcal W}\mathscr D_g
                       &=p(g)+\tau_0(g,\phi)-Z_0(g),\\
 \det_{\mathcal W}(1-z\mathscr E_g)
                       &=\det_{\mathcal W_+}(1+zL_g)
                                          \det_Q(1-zC_g).
 \end{aligned}
 \tag{MA-MO.855}
\]
The quotient determinant is the entire product
\(\prod_\rho(1-zG(\rho-1/2))^{m_\rho}\).
This triangular factorization is for \(\mathscr E_g\);
it does not assert that the raw \(\mathscr D_g\)
preserves \(\mathcal W_+\).

There is also an exact finite-rank correction formula
for the determinant of that raw matrix. In the actual
source-section coordinates let
\[
 \begin{gathered}
 V:\mathcal W\longrightarrow\mathbb C^2,\qquad V(u,c)=c,\\
 v_j=(-r_j,e_j),\qquad
 U_g:\mathbb C^2\longrightarrow\mathcal W,\qquad
 U_ge_j=C_gv_j ,
 \end{gathered}
 \tag{MA-MO.856}
\]
where \(r_j\) are (MA-MO.255) and \(e_j\) are the
standard polar coordinate vectors. Whenever
\(1+z\mathscr E_g\) is invertible,
\[
 \begin{aligned}
 \det_{\mathcal W}(1-z\mathscr D_g)
    ={}&\det_{\mathcal W_+}(1-zL_g)\det_Q(1+zC_g)\\
      &{}\cdot
       \det_{\mathbb C^2}\!
        \left(1-zV(1+z\mathscr E_g)^{-1}U_g\right).
 \end{aligned}
 \tag{MA-MO.857}
\]
The equality extends across the exceptional \(z\)'s
after removing the singularities of the product on
the right. No value of the two-by-two resolvent factor
alone is assigned at its poles.
For \(g(0)=0\), the raw trace is
\[
 \operatorname{Tr}_{\mathcal W}\mathscr D_g
       =\sum_p\sum_{k\ge1}
             \frac{\log p}{p^{k/2}}
                 \bigl[g(k\log p)+g(-k\log p)\bigr].
 \tag{MA-MO.858}
\]
The sums in (MA-MO.853) and (MA-MO.858) are finite.
For arbitrary \(g\), the trace statements are
(MA-MO.851) and (MA-MO.855), with the displayed
polar and nonsymmetric-cutoff terms retained.

\begin{proof}
\emph{Actual restrictions and the complete retained divisor.}
On the dense smooth ideal, the operator constructed in
MA-MO-027 and the theta composition in MA-MO-028
are the same original smoothed source action.
Both are bounded on \(J_\eta\). Their restrictions
therefore agree on all of \(J_\eta\), and
(MA-MO.803) gives its trace
\(\tau_{\eta,S}(g,\phi)\). Substitute this into
the already established actual triangular identity
(MA-MO.704) to prove (MA-MO.851) and the first
line of (MA-MO.852).

Every interior zero and every boundary jet is retained
under the explicit premises of this corollary.
MA-WP-001--003 therefore gives the full quotient
trace \(Z_\omega(g)\) on the actual Hilbert quotient.
On each finite zero block the diagonal value of
convolution is \(G(\rho-1/2)\), repeated according
to the full multiplicity. Nilpotent jet entries
do not alter a determinant. The trace-norm block
convergence proved in MA-WP, followed by Fredholm
determinant continuity, gives the stated infinite
products. They also converge locally uniformly
because \(\sum_\rho m_\rho|G(\rho-1/2)|<\infty\).
Thus neither a formally assigned divisor trace
nor a claim of simple zeros is used.

\emph{The rational relative operator on its correct invariant space.}
The bounded original negative projection is
\[
       P_-(u,c)=(u+c_0r_0+c_1r_1,0).
 \tag{MA-MO.859}
\]
Consequently \(F_g=U_gV\) has rank at most two.
The rank-one trace formula and the two actual
translation characters on the polar coordinates
give
\(\operatorname{Tr}F_g=G(-1/2)+G(1/2)=p(g)\).
This argument does not require the range of
\(1-P_-\) to be invariant under translation.
Since \(\mathscr D_g\) is trace class by
MA-MO-027, \(\mathscr E_g=F_g-\mathscr D_g\)
is trace class and agrees with the original
\(C_g(1-P_+)\) on the dense common source.

The actual source identity (MA-MO.447) is
\(p_+i_+=1+L\). Hence
\[
       C_g(1-P_+)i_+f=-i_+\mathcal D_gLf .
 \tag{MA-MO.860}
\]
The all-width return (MA-MO.829) supplies the
right side as \(-L_g\) on the completed
\(\mathcal W_+\). Density proves invariance
and the exact restriction, with trace
\(-\tau_0(g,\phi)\). The original \(C_gP_+\)
has range in the positive source image.
Passing to the actual quotient therefore
identifies the quotient action of
\(\mathscr E_g\) with \(C_g\). These are
exactly the source-level identities in
MA-MO-021; their narrow continuity is now
provided by MA-MO-027/028 instead of the
earlier strict half-plane margin.

In the orthogonal decomposition
\(\mathcal W=\mathcal W_+\oplus\mathcal W_+^\perp\),
\(\mathscr E_g\) is upper triangular with
diagonal blocks \(-L_g\) and the actual
quotient \(C_g\). Summing their absolutely
convergent diagonal traces and then taking
finite-rank compressions proves the trace
and determinant formulas for \(\mathscr E_g\).
Subtract its trace from \(\operatorname{Tr}F_g=p(g)\)
to obtain the raw trace in (MA-MO.855).

For the additional raw determinant formula,
\(\mathscr D_g=U_gV-\mathscr E_g\) gives
\[
  1-z\mathscr D_g
     =(1+z\mathscr E_g)
       \left(1-z(1+z\mathscr E_g)^{-1}U_gV\right).
 \tag{MA-MO.861}
\]
Fredholm determinant multiplicativity and the
finite-rank identity \(\det(1-AB)=\det(1-BA)\)
give a two-by-two determinant for the second
factor. Apply (MA-MO.855) with parameter
\(-z\) to the first factor. This proves
(MA-MO.857) on its stated resolvent set,
which contains a neighborhood of zero.
The finite-rank formula is meromorphic away
from the discrete exceptional set and agrees
there with the entire determinant on the
left; all singularities of its total product
are consequently removable. Indeed near an
exceptional point, the finite-dimensional
generalized eigenspace of the compact
\(\mathscr E_g\) gives the finite Laurent
principal part of its resolvent, while the
complementary resolvent is holomorphic.

\emph{Primitive prime evaluation and the value-zero restriction.}
For \(g(0)=0\), the unconditional primitive
explicit formulas and the actual local multiplier
calculation, in the already fixed conventions
of (MA-MO.484)--(MA-MO.486), are
\[
 \begin{aligned}
 Z_\eta(g)
    &=-\mathcal A_{\infty,\eta}(g)-B_\eta(g),\\
 \tau_{\eta,S}(g,\phi)
    &=-\mathcal A_{\infty,\eta}(g)-B_{\eta,S}(g),\\
 Z_0(g)&=p(g)-A_0(g)-B_0(g),\\
 \tau_0(g,\phi)&=-A_0(g).
 \end{aligned}
 \tag{MA-MO.862}
\]
Here \(A_0(g)\) abbreviates the same archimedean
term \(A_0(h)\) with \(h(e^x)=g(x)\) used
in MA-MO-021. The \(B\)'s are the full
primitive finite-prime distributions, and
\(B_{\eta,S}\) is their finite-place part
inside \(S\). In particular ramification of
the field alone does not delete an unramified
character factor. The negative test argument
has the conjugate character coefficient,
as written in (MA-MO.853).

The conductor and different monomials have
constant logarithmic derivatives and hence
pair with \(g(0)\); the compact residual is
also \(g(0)\int\psi\). Their disappearance
in (MA-MO.862) is exactly the stated
value-zero restriction. Their normalizations
remain present in (MA-MO.850) for general
tests. Substituting (MA-MO.862) into
(MA-MO.851) and (MA-MO.855) proves
(MA-MO.853) and (MA-MO.858), with their
positive prime signs.
\end{proof}

\paragraph{Scope and framework placement.}
The evaluated ideal trace is unconditional
in MA-MO-028. The full-matrix existence and
full quotient retention used here still
require precisely the divisor hypotheses
of MA-MO-027. In particular no new assertion
about the arithmetic zeros, and no
arbitrary no-margin sufficiency statement,
is a consequence of this corollary.

The rational Fredholm factorization for
the raw matrix retains the actual finite
rank-two factor in (MA-MO.857). The
rank-two outgoing obstruction is not
removed by relabeling the invariant
relative operator. In the inherited cote~1
framework, scans~079--080 and~111--112
through BR-TR-01/02, these are identities
of defined nuclear operators and their
actual invariant exact sequences. The
finite-length zero factors retain their
multiplicities without a diagonalizability
assumption; no new positivity or algebraic
weight claim is made.
